% GENERATED FILE. Edit canonical lessons, then run npm run generate:books.
% canonical-source-hash: fnv1a64:84b9224377cf7868
% canonical-lessons: TA-C120-kannati, TA-C120-coppu, TA-C120-vali, TA-C120-pena, TA-C120-kagitam

\chapter{Mirror, Soap, Bucket, Pen, Paper}
\label{ch:ta-mirror-soap-bucket-pen-paper}
\hlchaptermodality{120}

\begin{chapteropening}
\textbf{By the end of this chapter:} \emph{I can say a mirror, soap, a bucket, a pen and paper.}

Complete the last lesson of chapter 120: i can say a mirror, soap, a bucket, a pen and paper.
\end{chapteropening}

\section[kaṇṇāṭi]{\ta{கண்ணாடி} (kaṇṇāṭi) — a mirror; glasses}

\label{lesson:TA-C120-kannati}



\begin{quote}
Before the new one: say the Tamil for a pillow, then the Tamil for a blanket.
\end{quote}

\subsection*{You'll want to know: \ta{கண்ணாடி}}
\textbf{\ta{கண்ணாடி}} — \emph{kaṇṇāṭi} — \textquotedblleft{}a mirror; glasses\textquotedblright{}.

\begin{grammarlens}[title={What you've built}]
One more word for this chapter.
\end{grammarlens}

\subsection*{Your turn}
\begin{itemize}
\raggedright
  \item \emph{Say it:} \emph{kaṇṇāṭi}
  \item \emph{Say it:} \emph{kaṇṇāṭi}, once more
  \item \emph{Say it:} say \emph{kaṇṇāṭi}
  \item \emph{Recall:} say the Tamil for a pillow, then the Tamil for a blanket, then say \emph{kaṇṇāṭi} again
\end{itemize}

\begin{tcolorbox}[breakable,colback=teal!4,colframe=teal!35!black,title={Before you move on}]
What does \ta{கண்ணாடி} mean? (\textquotedblleft{}A mirror; glasses\textquotedblright{}.) Say it once more.
\end{tcolorbox}
\section[cōppu]{\ta{சோப்பு} (cōppu) — soap}

\label{lesson:TA-C120-coppu}



\begin{quote}
Before the new one: say the Tamil for a blanket, then the Tamil for a mirror; glasses.
\end{quote}

\subsection*{You'll want to know: \ta{சோப்பு}}
\textbf{\ta{சோப்பு}} — \emph{cōppu} — \textquotedblleft{}soap\textquotedblright{}.

\begin{grammarlens}[title={What you've built}]
One more word for this chapter.
\end{grammarlens}

\subsection*{Your turn}
\begin{itemize}
\raggedright
  \item \emph{Say it:} \emph{cōppu}
  \item \emph{Say it:} \emph{cōppu}, once more
  \item \emph{Say it:} say \emph{cōppu}
  \item \emph{Recall:} say the Tamil for a blanket, then the Tamil for a mirror; glasses, then say \emph{cōppu} again
\end{itemize}

\begin{tcolorbox}[breakable,colback=teal!4,colframe=teal!35!black,title={Before you move on}]
What does \ta{சோப்பு} mean? (\textquotedblleft{}Soap\textquotedblright{}.) Say it once more.
\end{tcolorbox}
\section[vāḷi]{\ta{வாளி} (vāḷi) — a bucket}

\label{lesson:TA-C120-vali}



\begin{quote}
Before the new one: say the Tamil for a mirror; glasses, then the Tamil for soap.
\end{quote}

\subsection*{You'll want to know: \ta{வாளி}}
\textbf{\ta{வாளி}} — \emph{vāḷi} — \textquotedblleft{}a bucket\textquotedblright{}.

\begin{grammarlens}[title={What you've built}]
One more word for this chapter.
\end{grammarlens}

\subsection*{Your turn}
\begin{itemize}
\raggedright
  \item \emph{Say it:} \emph{vāḷi}
  \item \emph{Say it:} \emph{vāḷi}, once more
  \item \emph{Say it:} say \emph{vāḷi}
  \item \emph{Recall:} say the Tamil for a mirror; glasses, then the Tamil for soap, then say \emph{vāḷi} again
\end{itemize}

\begin{tcolorbox}[breakable,colback=teal!4,colframe=teal!35!black,title={Before you move on}]
What does \ta{வாளி} mean? (\textquotedblleft{}A bucket\textquotedblright{}.) Say it once more.
\end{tcolorbox}
\section[pēṉā]{\ta{பேனா} (pēṉā) — a pen}

\label{lesson:TA-C120-pena}



\begin{quote}
Before the new one: say the Tamil for soap, then the Tamil for a bucket.

Say \textbf{\ta{தமிழ்}}, curling the tongue back for \textbf{\ta{ழ}}, then \textbf{\ta{எழுது}}.
\end{quote}

\subsection*{You'll want to know: \ta{பேனா}}
\textbf{\ta{பேனா}} — \emph{pēṉā} — \textquotedblleft{}a pen\textquotedblright{}.

\begin{grammarlens}[title={What you've built}]
One more word for this chapter.
\end{grammarlens}

\subsection*{Your turn}
\begin{itemize}
\raggedright
  \item \emph{Say it:} \emph{pēṉā}
  \item \emph{Say it:} \emph{pēṉā}, once more
  \item \emph{Say it:} say \emph{pēṉā}
  \item \emph{Recall:} say the Tamil for soap, then the Tamil for a bucket, then say \emph{pēṉā} again
\end{itemize}

\begin{tcolorbox}[breakable,colback=teal!4,colframe=teal!35!black,title={Before you move on}]
What does \ta{பேனா} mean? (\textquotedblleft{}A pen\textquotedblright{}.) Say it once more.
\end{tcolorbox}
\section[kāgitam]{\ta{காகிதம்} (kāgitam) — paper}

\label{lesson:TA-C120-kagitam}



\begin{quote}
Before the new one: say the Tamil for a bucket, then the Tamil for a pen.

Read the lines aloud as one run: \textbf{\ta{வணக்கம்}. \ta{இது} \ta{என்} \ta{வீடு}. \ta{அம்மா} \ta{இங்கே}.}
\end{quote}

\subsection*{You'll want to know: \ta{காகிதம்}}
\textbf{\ta{காகிதம்}} — \emph{kāgitam} — \textquotedblleft{}paper\textquotedblright{}.

\begin{grammarlens}[title={What you've built}]
One more word for this chapter.
\end{grammarlens}

\subsection*{Your turn}
\begin{itemize}
\raggedright
  \item \emph{Say it:} \emph{kāgitam}
  \item \emph{Say it:} \emph{kāgitam}, once more
  \item \emph{Say it:} say \emph{kāgitam}
  \item \emph{Recall:} say the Tamil for a bucket, then the Tamil for a pen, then say \emph{kāgitam} again
\end{itemize}

\begin{tcolorbox}[breakable,colback=teal!4,colframe=teal!35!black,title={Before you move on}]
What does \ta{காகிதம்} mean? (\textquotedblleft{}Paper\textquotedblright{}.) Say it once more.
\end{tcolorbox}
